\subsection{An odd cyclic primitive factor in every genus and characteristic}

\begin{claim}{A.C2.20}{Proved}
Let $K$ be finitely generated over $\F_p$, and let
$C/K$ be smooth proper geometrically connected of
genus $g\ge2$. Suppose that geometrically there is
a connected cyclic \'etale cover $D\to C$ of odd
degree $n>1$ prime to $p$. For a deck generator
$\sigma$, put
\[
 P_n=\operatorname{im}(A)\subset J_D,\qquad
                 A=\prod_{\ell\mid n}(1-\sigma^{n/\ell}),
                                                               \tag{GP.1}
\]
where the product runs over distinct prime divisors.
If $f_p(P_n)=0$, then the geometrically prime-to-$p$
point-section map of $C$ is injective.
The conclusion holds after every finite extension,
for forms satisfying this geometric hypothesis,
and for hyperbolic open subcurves.

The image in (GP.1) is the primitive conductor-$n$
factor up to isogeny: its scalar on character $i$
is nonzero exactly when $\gcd(i,n)=1$.
Its isogeny class is independent of the generator.
In genus two and characteristic two, A.C2.5
identifies this hypothesis with the primitive
doubling-orbit condition in A.C2.19.
No theta characterization is asserted here in
higher genus.
\end{claim}
\begin{proof}
Suppose that distinct $P,Q\in C(K)$ give equal
point-section classes. After a finite extension,
define $D$, its action, and a chosen lift $P'$.
By A.3 there is a lift $Q'$ with
$[P'-Q']\in J_D[p^\infty]$.
The image of this class under $A$ belongs to
$P_n[p^\infty]$, which is zero by hypothesis.

Write $r=\omega(n)$, $d=2^r$, and
$a_S=\prod_{\ell\in S}\sigma^{n/\ell}$ for a
subset $S$ of the distinct primes dividing $n$.
Thus
\[
                 \sum_S(-1)^{|S|}(a_SP'-a_SQ')            \tag{GP.2}
\]
is principal. The prime-order factors have
pairwise coprime orders, so the $a_S$ are distinct.
The two deck orbits are free and disjoint.
Consequently (GP.2) has $d$ distinct zeros and
$d$ distinct poles, each simple. Its rational
function $f$ is separable of degree $d$.

For every odd $n>1$ one has
\[
                              n>(d-1)^2.                 \tag{GP.3}
\]
For $r=1$, use $n\ge3>1$; for $r=2$, use
$n\ge15>9$; for $r\ge3$, use
$n\ge3\cdot5^{r-1}>4^r>(2^r-1)^2$.
Let
\[
 M=\overline K(\gamma f:\gamma\in G),\qquad
                    e=[\overline K(D):M],
\]
and let $Z$ be its smooth proper curve.
The controlled-field induction proved in A.C2.19,
which uses only Castelnuovo--Severi, gives
\[
                       e\mid d,\qquad
                       g(Z)\le(d/e-1)^2.                 \tag{GP.4}
\]
It is valid in every characteristic
\citep[Proposition 2.1, Lemma 2.2 and Corollary 2.3,
author p.~3]{A_PrimPoonen2007}.
Since $g(D)=n(g-1)+1\ge n+1$, (GP.3) excludes
$e=1$. Thus $e>1$ is a power of two.

The field $M$ is $G$-stable. Its induced action
is faithful, since the kernel order divides both
$e$ and the odd number $n$. It is free: any
odd-prime-order stabilizer would partition the
degree-$e$ fibre divisor in $D$ into free orbits
with constant multiplicity, forcing that prime
to divide $e$.
Hence $g(Z)-1=n(g(Z/G)-1)$. Together with
(GP.3)--(GP.4), this forces $g(Z)=1$.

Factor $f=h\circ v$ through $v:D\to Z$.
The fibre above a zero of $h$ consists of
exactly $e$ unramified points, since the zero
divisor (GP.2) is reduced. Each of the two
original deck orbits maps injectively to $Z$,
by its free $G$-action. The whole zero fibre
lies in these two orbits, so $e\le2$. Thus
$e=2$. The resulting separable quadratic
involution $\tau$ commutes with $G$, since
$M$ is $G$-stable. For $r=1$, (GP.4) would
already contradict $g(Z)=1$; the following
argument covers every hypothetical quadratic
quotient obtained this way.

Put $E_0=Z/G$. The commuting actions give
$G\times\langle\tau\rangle$ on $D$, and
\[
 \overline K(E_0)
       =\overline K(D)^{G\times\langle\tau\rangle}.
\]
Thus $E_0$ has genus one and $C=D/G\to E_0$
is separable quadratic. It is ramified,
because $g(C)\ge2$. The map $Z\to E_0$ is
an odd \'etale isogeny after choosing origins.

For every $m\mid n$ put
$H_m=\langle\sigma^m\rangle$,
$D_m=D/H_m$, and $Z_m=Z/H_m$. The square
\[
 \begin{array}{ccc}
 D_m&\longrightarrow&Z_m\\
 \downarrow&&\downarrow\\
 C&\longrightarrow&E_0
 \end{array}                                              \tag{GP.5}
\]
has quadratic horizontal maps and degree-$m$
\'etale vertical maps. It is Cartesian.
Indeed $W=C\times_{E_0}Z_m$ has $W\to C$
\'etale, as the base change of $Z_m\to E_0$,
and hence $W$ is smooth. The extensions of
degrees two and $m$ over $\overline K(E_0)$
are linearly disjoint. Thus $W$ is integral,
and their compositum, of degree $2m$, is
$\overline K(D_m)$. This identifies $W$ with
$D_m$.

Suppose first that $p=2$. Let $r_1>0$ be
the number of geometric ramification points
of $C\to E_0$. By (GP.5), the corresponding
number for $D_m\to Z_m$ is $m r_1$:
every branch point has $m$ distinct preimages
under the \'etale isogeny and the same local
quadratic extension after base change.
All the genus-one curves $Z_m$ have the
same two-rank $\epsilon$, by isogeny.
The ramified Deuring--Shafarevich formula gives
\[
                       f_2(D_m)=m r_1+2\epsilon-1
                                                               \tag{GP.6}
\]
\citep[Corollary 1.8, pp.~35--36]{A_PrimCrew1984}.
Here $r_1$ is positive by Riemann--Hurwitz,
not an expression whose sign is left to
be determined.

Let $b_j$ be the rank of the exact
conductor-$j$ part of the geometric \'etale
two-adic Tate representation of $J_D$.
The odd-order averaging idempotent, together
with pullback and norm, identifies its
$H_m$-invariants with the Tate module of
$J_{D_m}$. Therefore
\[
                          f_2(D_m)=\sum_{j\mid m}b_j.
\]
The character scalars in (GP.1) identify
$b_n=f_2(P_n)$. M\"obius inversion yields
\[
 \begin{aligned}
 f_2(P_n)
 &=\sum_{m\mid n}\mu(n/m)f_2(D_m)\\
 &=r_1\sum_{m\mid n}\mu(n/m)m
   +(2\epsilon-1)\sum_{m\mid n}\mu(n/m)\\
 &=r_1\varphi(n)>0,
 \end{aligned}                                             \tag{GP.7}
\]
contradicting the hypothesis. No
one-dimensional character-space assumption
was used.

If $p$ is odd, the derived quadratic map
$C\to E_0$ is instead a geometrically
prime-to-$p$ Galois cover of a genus-one
curve. A.29l already gives point-section
injectivity for exactly this class, including
its finite-extension and finitely generated
form scopes. It excludes the assumed
collision. No Deuring--Shafarevich formula
for a group of order prime to $p$ is invoked.

The cover and quotient data descend after
finite extension, equality of section classes
restricts to that extension, and A.3 supplies
the required lifted $p$-primary difference.
The geometric contradiction therefore
establishes the assertion over $K$ and over
every finite extension. Forms are included
by the same reasoning, and equality on an
open curve maps to equality on its
compactification.

This is a sufficient geometric condition,
not an existence theorem for primitive
factors of zero $p$-rank. Nonordinariness
of a cover is weaker and does not imply it.
The certificate
\path{code/BC_general_primitive_prym_checks.py}
and its JSON output verify the M\"obius
identity symbolically at 250 odd orders and
in 56 numerical rank examples. The proof
above gives the general identity and its
geometric meaning.
\end{proof}
